
Data curation for foundation model training typically combines quality filtering and diversity sampling, yet optimal ordering of these operations remains underexplored. Through systematic experiments on GPT-2 Small across dataset scales (10K--10M tokens), we test whether quality-first (QD) or diversity-first (DQ) curation produces superior downstream performance. Our findings reveal that quality-first ordering universally outperforms diversity-first by 1.5--5.0 percentage points (pp) across all tested scales (Cohen's d=0.76--2.48). We quantify quality saturation, showing that quality-only improvement decreases from 11.0pp at 10K tokens to 3.3pp at 10M tokens (slope $-2.62pp/log-scale$, p=0.008, $R^{2}=0.93)$. Diversity persistence is modeled from literature showing sustained effectiveness at large scale. Contrary to predictions of scale-dependent reversal (diversity-first dominating at large scale), quality-first curation maintains superiority even when diversity coefficient exceeds quality coefficient at 10M tokens. This suggests a quality-gated diversity mechanism: diversity sampling appears most effective when applied to quality-filtered subsets rather than raw corpora. For practitioners training 100M--1B parameter models on web corpora at 10K--10M scale, we provide prescriptive guidance: apply quality filtering before diversity sampling.
